% !TEX program = xelatex
\documentclass[10pt,letterpaper,twocolumn]{article}

% === GEOMETRY ===
\usepackage[top=0.75in, bottom=0.75in, left=0.75in, right=0.75in]{geometry}

% === FONTS ===
\usepackage{fontspec}
\usepackage{unicode-math}
\usepackage{xeCJK}
\setCJKmainfont{Droid Sans Fallback}
\setmainfont{Latin Modern Roman}
\setsansfont{Latin Modern Sans}
\setmonofont{Latin Modern Mono}[Scale=0.85]
\newfontfamily\acbold{ywft-processing-bold}[
  Path=../../system/public/type/webfonts/,
  Extension=.ttf
]
\newfontfamily\aclight{ywft-processing-light}[
  Path=../../system/public/type/webfonts/,
  Extension=.ttf
]

% === PACKAGES ===
\usepackage{xcolor}
\usepackage{titlesec}
\usepackage{enumitem}
\usepackage{booktabs}
\usepackage{tabularx}
\usepackage{fancyhdr}
\usepackage{hyperref}
\usepackage{graphicx}
\usepackage{ragged2e}
\usepackage{listings}
\usepackage{natbib}
\usepackage[colorspec=0.92]{draftwatermark}

% === COLORS ===
\definecolor{acpink}{RGB}{180,72,135}
\definecolor{acpurple}{RGB}{120,80,180}
\definecolor{acdark}{RGB}{64,56,74}
\definecolor{acgray}{RGB}{119,119,119}
\definecolor{draftcolor}{RGB}{180,72,135}

% === DRAFT WATERMARK ===
\DraftwatermarkOptions{
  text=WORKING DRAFT,
  fontsize=3cm,
  color=draftcolor!18,
  angle=45,
  pos={0.5\paperwidth, 0.5\paperheight}
}

% === JS SYNTAX COLORS ===
\definecolor{jskw}{RGB}{119,51,170}
\definecolor{jsfn}{RGB}{0,136,170}
\definecolor{jsstr}{RGB}{170,120,0}
\definecolor{jsnum}{RGB}{204,0,102}
\definecolor{jscmt}{RGB}{102,102,102}

% === PROCESSING SYNTAX COLORS ===
\definecolor{pjkw}{RGB}{204,102,0}
\definecolor{pjfn}{RGB}{0,102,153}
\definecolor{pjcmt}{RGB}{102,102,102}

% === HYPERREF ===
\hypersetup{
  colorlinks=true,
  linkcolor=acpurple,
  urlcolor=acpurple,
  citecolor=acpurple,
  pdfauthor={@jeffrey},
  pdftitle={从setup()到boot()：Processing在作品API核心中的延续},
}

% === SECTION FORMATTING ===
\titleformat{\section}
  {\normalfont\bfseries\normalsize\uppercase}
  {\thesection.}
  {0.5em}
  {}
\titlespacing{\section}{0pt}{1.2em}{0.3em}

\titleformat{\subsection}
  {\normalfont\bfseries\small}
  {\thesubsection}
  {0.5em}
  {}
\titlespacing{\subsection}{0pt}{0.8em}{0.2em}

\titleformat{\subsubsection}
  {\normalfont\itshape\small}
  {\thesubsubsection}
  {0.5em}
  {}
\titlespacing{\subsubsection}{0pt}{0.6em}{0.1em}

% === HEADER/FOOTER ===
\pagestyle{fancy}
\fancyhf{}
\renewcommand{\headrulewidth}{0pt}
\fancyhead[C]{\footnotesize\color{draftcolor}\textit{工作草稿 --- 请勿引用}}
\fancyfoot[C]{\footnotesize\thepage}

% === LIST SETTINGS ===
\setlist[itemize]{nosep, leftmargin=1.2em, itemsep=0.1em}
\setlist[enumerate]{nosep, leftmargin=1.2em}

% === COLUMN SEPARATION ===
\setlength{\columnsep}{1.8em}

% === PARAGRAPH SETTINGS ===
\setlength{\parindent}{1em}
\setlength{\parskip}{0.3em}

% Hyphenation for narrow two-column layout
\tolerance=800
\emergencystretch=1em
\hyphenpenalty=50

% === LISTINGS ===
\lstdefinelanguage{processing}{
  morekeywords=[1]{void,int,float,boolean,color,String,class,new,if,else,for,while,return,public,private,extends,import},
  morekeywords=[2]{setup,draw,mousePressed,mouseDragged,keyPressed,size,background,stroke,noStroke,fill,noFill,line,rect,ellipse,point,triangle,beginShape,endShape,vertex,translate,rotate,scale,pushMatrix,popMatrix,frameRate,width,height,mouseX,mouseY,pmouseX,pmouseY,key,keyCode,mouseButton,frameCount},
  sensitive=true,
  morecomment=[l]{//},
  morecomment=[s]{/*}{*/},
  morestring=[b]",
}

\lstdefinelanguage{p5js}{
  morekeywords=[1]{function,let,const,var,if,else,for,while,return,new,class,export},
  morekeywords=[2]{setup,draw,mousePressed,mouseDragged,keyPressed,createCanvas,background,stroke,noStroke,fill,noFill,line,rect,ellipse,point,triangle,beginShape,endShape,vertex,translate,rotate,scale,push,pop,frameRate,width,height,mouseX,mouseY,pmouseX,pmouseY,key,keyCode,mouseButton,frameCount,createGraphics,random,noise,map,constrain,dist,lerp},
  sensitive=true,
  morecomment=[l]{//},
  morestring=[b]",
  morestring=[b]',
  morestring=[b]`,
}

\lstdefinelanguage{acjs}{
  morekeywords=[1]{function,export,const,let,var,return,if,else,new,async,await,import,from,for},
  morekeywords=[2]{wipe,ink,line,box,circle,write,screen,params,colon,jump,send,store,net,sound,speaker,pen,event,boot,paint,act,sim,leave,paste,plot,poly,flood,form,cursor,handle,num,geo,ui},
  sensitive=true,
  morecomment=[l]{//},
  morestring=[b]",
  morestring=[b]',
  morestring=[b]`,
}

\lstdefinestyle{processingstyle}{
  language=processing,
  keywordstyle=[1]\color{pjkw}\bfseries,
  keywordstyle=[2]\color{pjfn}\bfseries,
  commentstyle=\color{pjcmt}\itshape,
  stringstyle=\color{jsstr},
}

\lstdefinestyle{p5style}{
  language=p5js,
  keywordstyle=[1]\color{jskw}\bfseries,
  keywordstyle=[2]\color{pjfn}\bfseries,
  commentstyle=\color{jscmt}\itshape,
  stringstyle=\color{jsstr},
}

\lstdefinestyle{acjsstyle}{
  language=acjs,
  keywordstyle=[1]\color{jskw}\bfseries,
  keywordstyle=[2]\color{jsfn}\bfseries,
  commentstyle=\color{jscmt}\itshape,
  stringstyle=\color{jsstr},
}

\lstset{
  basicstyle=\ttfamily\small,
  breaklines=true,
  frame=single,
  rulecolor=\color{acgray!30},
  backgroundcolor=\color{acgray!5},
  xleftmargin=0.5em,
  xrightmargin=0.5em,
  aboveskip=0.5em,
  belowskip=0.5em,
  numbers=none,
  tabsize=2,
}

\newcommand{\acdot}{{\color{acpink}.}}
\newcommand{\ac}{\textsc{Aesthetic.Computer}}

\begin{document}

% ============ TITLE BLOCK ============

\twocolumn[{%
\begin{center}
{\acbold\fontsize{22pt}{26pt}\selectfont\color{acdark} 从 \texttt{setup()} 到 \texttt{boot()}}\par
\vspace{0.2em}
{\aclight\fontsize{11pt}{13pt}\selectfont\color{acpink} Processing在作品API核心中的延续}\par
\vspace{0.6em}
{\normalsize @jeffrey}\par
{\small\color{acgray} Aesthetic.Computer}\par
{\small\color{acgray} ORCID: \href{https://orcid.org/0009-0007-4460-4913}{0009-0007-4460-4913}}\par
\vspace{0.3em}
{\small\color{acpurple} \url{https://aesthetic.computer}}\par
\vspace{0.6em}
\rule{\textwidth}{1.5pt}
\vspace{0.5em}
\end{center}

\begin{center}
{\small\color{draftcolor}\textbf{[ 工作草稿 --- 请勿引用 ]}}
\end{center}
\vspace{0.3em}

\begin{quote}
\small\noindent\textbf{摘要。}
每个创意编程平台都定义了一套核心API——一组使程序感知世界并在其上绘制的基本函数。Processing在2001年以\texttt{setup()}/\texttt{draw()}加全局绘图原语（\texttt{background()}、\texttt{stroke()}、\texttt{ellipse()}）确立了一种模式，这一模式影响了二十年的创意编程工具。p5.js在2014年将此模型引入浏览器，在适应JavaScript和DOM的同时保留了API接口。\ac{}（AC）始于2021年，其核心延续了Processing的思想，同时在多个结构层面进行了扩展：它将两函数生命周期扩展为五个（\texttt{boot}、\texttt{paint}、\texttt{act}、\texttt{sim}、\texttt{leave}），以解构式API注入取代全局状态，以不同频率分离模拟与渲染，并使每个程序无需构建步骤即可通过URL寻址。本文追溯Processing的理念如何存在于AC作品API中，比较了Design By Numbers、Processing、p5.js和AC在生命周期模型、绘图原语、输入处理和状态管理策略方面的差异。我们论证，AC的API设计——以Processing的即时模式图形模型为基础——将\emph{速写本}隐喻（编写、运行、丢弃）扩展为\emph{乐器}隐喻（启动、演奏、练习、回归）。
\end{quote}
\vspace{0.5em}
}]

% ============ 1. INTRODUCTION ============

\section{引言}

创意编程API的历史，是关于决定程序员应首先说什么的历史。在Design By Numbers~\citep{maeda2001dbn}中，第一条语句是坐标和灰度值：\texttt{Paper 50}。在Processing~\citep{reas2003processing}中，是\texttt{setup()}内的\texttt{size(200, 200)}。在p5.js~\citep{mccarthy2015p5js}中，是\texttt{createCanvas(400, 400)}。在\ac{}中，程序员完全无需声明画布——运行时将\texttt{screen.width}和\texttt{screen.height}作为既定事实提供，第一条有意义的语句通常是\texttt{wipe()}，清除屏幕以开始绘制。

这些不是任意的差异。每一个都反映了平台假设与程序员必须指定之间的决策，这些决策累积后形成了作者与机器之间根本不同的关系。本文从Processing到p5.js再到AC追溯API谱系，将每次转变视为一系列设计决策：保留了什么，改变了什么，以及改变意味着什么。

本文的贡献有三方面：（1）对Processing、p5.js和AC的生命周期模型、绘图API、输入系统和状态管理策略进行详细技术比较；（2）分析AC偏离Processing模型背后的设计理念；（3）论证AC作品API体现了从速写本隐喻到我们所称的\emph{乐器隐喻}的转变——一种程序员与程序之间关系持续进行、具有表演性、基于练习而非探索性和一次性的模型。

% ============ 2. THE PROCESSING MODEL ============

\section{Processing模型（2001）}
\label{sec:processing}

Processing~\citep{reas2007processing}被设计为面向视觉艺术家的"软件速写本"。其API以两个生命周期函数和一组全局绘图原语为中心。

\subsection{生命周期：\texttt{setup()}和\texttt{draw()}}

\begin{lstlisting}[style=processingstyle,caption={最小Processing草图。}]
void setup() {
  size(400, 400);
  background(0);
}

void draw() {
  stroke(255);
  ellipse(mouseX, mouseY, 20, 20);
}
\end{lstlisting}

\texttt{setup()}运行一次；\texttt{draw()}以默认60fps的频率每帧运行。这种两函数模型是Processing最具影响力的设计决策。它将动画的认知开销从管理渲染循环降低到填写两个空白："发生一次什么？"和"每帧发生什么？"

该模型刻意保持最小化。没有显式的\texttt{input()}函数——鼠标和键盘状态以自动更新的全局变量（\texttt{mouseX}、\texttt{mouseY}、\texttt{key}、\texttt{keyCode}）形式可用。事件回调（\texttt{mousePressed()}、\texttt{keyPressed()}）存在但是可选的补充，而非必需的入口点。

\subsection{绘图原语}

Processing的绘图API采用继承自PostScript和OpenGL的\emph{有状态管道}模型：状态设置调用（\texttt{stroke()}、\texttt{fill()}、\texttt{strokeWeight()}）修改持久的图形上下文，形状调用（\texttt{ellipse()}、\texttt{rect()}、\texttt{line()}）使用该上下文进行渲染。上下文在帧间持续存在，除非显式重置。

\begin{lstlisting}[style=processingstyle,caption={有状态绘图上下文。}]
void draw() {
  fill(255, 0, 0);     // Set fill: red
  stroke(0);            // Set stroke: black
  strokeWeight(3);      // Set weight: 3px
  ellipse(100, 100, 50, 50);  // Uses above
  rect(200, 100, 50, 50);     // Also uses above
  // State persists into next frame
}
\end{lstlisting}

这种状态持久性既强大（统一样式所需代码更少）又容易出错（被遗忘的状态在帧间泄漏）。Processing通过\texttt{pushStyle()}/\texttt{popStyle()}和\texttt{pushMatrix()}/\texttt{popMatrix()}来解决这个问题，但这些是高级工具。

\subsection{全局作用域}

所有Processing程序共享一个全局命名空间。在\texttt{setup()}和\texttt{draw()}外部声明的变量在任何地方都可访问。绘图函数（\texttt{line()}、\texttt{ellipse()}、\texttt{background()}）是全局的。输入变量（\texttt{mouseX}、\texttt{mouseY}）是全局的。这消除了理解作用域、导入或依赖注入的需要——这对初学者至关重要——但使程序无法组合或隔离。

% ============ 3. THE P5.JS TRANSITION ============

\section{p5.js的过渡（2014）}
\label{sec:p5js}

p5.js~\citep{mccarthy2015p5js}是Processing的JavaScript重新实现。其主要贡献是将Processing模型带入浏览器，但移植需要若干适配。

\subsection{生命周期的保留}

\begin{lstlisting}[style=p5style,caption={最小p5.js草图。}]
function setup() {
  createCanvas(400, 400);
}

function draw() {
  background(220);
  ellipse(mouseX, mouseY, 50, 50);
}
\end{lstlisting}

\texttt{setup()}/\texttt{draw()}模型被完全保留。认知模型相同：两个函数、全局绘图原语、自动动画循环。这种连续性是刻意的——p5.js旨在成为Web端的Processing，而非一种新语言。

\subsection{画布创建}

最明显的变化是\texttt{createCanvas()}替代了\texttt{size()}。这不是表面的改变：在Processing中，\texttt{size()}配置应用程序窗口；在p5.js中，\texttt{createCanvas()}在DOM内创建一个HTML \texttt{<canvas>}元素。草图必须与现有页面协商，而非独占整个显示区域。

\subsection{实例模式}

p5.js引入了一个重要的全局命名空间逃逸出口：\emph{实例模式}。

\begin{lstlisting}[style=p5style,caption={p5.js实例模式。}]
const sketch = (p) => {
  p.setup = () => {
    p.createCanvas(400, 400);
  };
  p.draw = () => {
    p.background(220);
    p.ellipse(p.mouseX, p.mouseY, 50, 50);
  };
};
new p5(sketch);
\end{lstlisting}

实例模式将所有API函数包装在命名空间（\texttt{p.}）后面，允许一个页面上有多个草图并避免全局污染。然而，大多数p5.js代码以全局模式编写，实例模式通常只在将草图嵌入更大应用程序时才遇到。前缀记法（\texttt{p.ellipse}对比\texttt{ellipse}）增加了冗余性，降低了使Processing吸引人的"直接编写"品质。

\subsection{改变了什么，没改变什么}

p5.js保留了：两函数生命周期、全局绘图原语、有状态图形上下文、基于帧的动画以及"草图"身份。它适配了：画布创建适配DOM、渲染器选择（2D/WebGL）、事件处理适配浏览器事件，并添加了用于组合的实例模式。它没有解决：模拟与渲染的混淆、结构化输入处理的缺失，以及单文件隔离问题。

% ============ 4. THE AC PIECE API ============

\section{AC作品API（2021--）}
\label{sec:ac}

\ac{}的作品API~\citep{scudder2026ac}建立在Processing的核心理念之上——即时模式图形、单文件程序、生命周期函数——并围绕五个关注点而非两个进行扩展。

\subsection{五个生命周期函数}

\begin{lstlisting}[style=acjsstyle,caption={最小AC作品。}]
function boot({ screen, params }) {
  // Runs once on load
}

function paint({ wipe, ink, line, screen }) {
  wipe("navy");
  ink("pink").circle(
    screen.width / 2,
    screen.height / 2,
    50
  );
}

function act({ event: e }) {
  if (e.is("keyboard:down:space")) {
    // Handle spacebar
  }
}

function sim() {
  // Update state at 120fps
}

export { boot, paint, act, sim };
\end{lstlisting}

五个函数是：

\begin{itemize}
  \item \texttt{boot(\$api)} --- 作品加载时运行一次。扩展了Processing的\texttt{setup()}，由运行时提供屏幕而非由程序员配置。
  \item \texttt{paint(\$api)} --- 需要渲染的每一帧运行。延续了\texttt{draw()}的即时模式模型，但纯粹是视觉的：按惯例不包含逻辑、不修改状态。
  \item \texttt{act(\$api)} --- 每个输入事件调用一次。将Processing的事件回调（\texttt{mousePressed()}、\texttt{keyPressed()}）整合为单一的统一入口点。
  \item \texttt{sim(\$api)} --- 以固定120Hz运行，可能在每个渲染帧期间运行多次。在Processing中没有等价物；最近的类比是固定时间步长游戏循环~\citep{nystrom2014game}。
  \item \texttt{leave(\$api)} --- 退出时的清理。在Processing中没有等价物，草图只是简单地终止。
\end{itemize}

每个函数都是可选的。只导出\texttt{paint}的作品也是一个有效的、可运行的程序。

\subsection{通过解构的API注入}

对Processing模型最具结构性意义的扩展是\textbf{全局作用域中不存在任何API函数}。作品使用的每个函数都作为参数接收：

\begin{lstlisting}[style=acjsstyle,caption={AC中的API解构。}]
function paint({ wipe, ink, line, box,
                 circle, screen, num }) {
  wipe(0, 0, 30);
  ink(255, 100, 180);
  const cx = num.lerp(0, screen.width, 0.5);
  circle(cx, screen.height / 2, 40);
}
\end{lstlisting}

这种设计有几个后果：

\begin{enumerate}
  \item \textbf{无全局污染。}模块级变量是作品状态；API函数是参数。不存在可能冲突的环境命名空间。
  \item \textbf{自文档化的导入。}每个函数顶部的解构模式精确声明了该函数使用的API接口，充当内联依赖清单。
  \item \textbf{可组合性。}多个作品可以在同一JavaScript上下文中共存而不会产生API冲突，实现了平台的作品切换和嵌入机制。
  \item \textbf{可发现性。}编辑器中的自动补全对解构的参数对象生效，提供了探索API的自然方式。
\end{enumerate}

这是p5.js实例模式模式推到其结论：程序员不是在每次调用前加\texttt{p.}前缀，而是在函数签名处解构出恰好需要的函数，只需一次。

\subsection{模拟与渲染的分离}
\label{sec:simsep}

Processing的\texttt{draw()}混淆了两个关注点：更新状态和渲染像素。这对简单草图有效，但随着复杂度增长会产生问题：物理模拟绑定到帧率、输入延迟耦合到渲染成本，以及难以分离"发生了什么"和"看起来像什么"。

AC显式地分离了这些：

\begin{itemize}
  \item \texttt{sim()}以固定120Hz运行，与显示刷新率无关。在60Hz显示器上，\texttt{sim()}在每次\texttt{paint()}调用中运行两次。在144Hz显示器上，比率相应调整。
  \item \texttt{paint()}以显示器刷新率运行（上限约165fps），仅负责渲染当前状态。
  \item \texttt{act()}在事件到达时立即运行，每个输入事件一次。
\end{itemize}

这种三向分离反映了现代游戏引擎的架构~\citep{nystrom2014game}：用于确定性逻辑的固定更新频率、用于平滑显示的可变渲染频率，以及用于输入的事件队列。区别在于AC将此架构作为第一等API暴露，而非隐藏在框架之后。

\subsection{统一的事件处理}

Processing将输入处理分散在全局变量和可选回调中：

\begin{lstlisting}[style=processingstyle]
// Processing: scattered input model
void draw() {
  if (mousePressed) { /* poll state */ }
}
void mousePressed() { /* callback */ }
void keyPressed()   { /* callback */ }
void mouseDragged() { /* callback */ }
\end{lstlisting}

AC通过基于字符串的事件匹配系统将所有输入整合到\texttt{act()}中：

\begin{lstlisting}[style=acjsstyle]
function act({ event: e, pen }) {
  if (e.is("touch"))              { }
  if (e.is("draw"))               { }
  if (e.is("lift"))               { }
  if (e.is("keyboard:down:a"))    { }
  if (e.is("keyboard:down:space")){ }
  if (e.is("gamepad:a"))          { }
  if (e.is("reframed"))           { }
}
\end{lstlisting}

\texttt{event.is()}模式通过单一函数提供了跨输入模态（触控、鼠标、键盘、手柄、MIDI、手部追踪）的统一接口。事件字符串是分层的（\texttt{keyboard:down:a}），可以在任何特异性级别进行匹配。这用单一的、可过滤的流替代了Processing回调的组合爆炸。

\subsection{绘图原语：默认无状态}

Processing和p5.js使用有状态的图形上下文，其中\texttt{fill()}、\texttt{stroke()}和变换调用持续到显式更改。AC反转了这一点：\texttt{ink()}函数为\emph{紧随其后的}绘图调用设置颜色，并返回API对象以支持链式调用：

\begin{lstlisting}[style=acjsstyle]
function paint({ wipe, ink, line, circle }) {
  wipe("black");
  ink("red").circle(50, 50, 20);
  ink("blue").line(0, 0, 100, 100);
  // No state leaks between calls
}
\end{lstlisting}

\texttt{ink()}的返回值链式调用模式使颜色-形状对在视觉上是原子的：每个绘图操作都是自包含的表达式，而非一系列上下文突变。这消除了一类错误：来自上一帧（或上一个函数）的被遗忘的\texttt{fill()}或\texttt{stroke()}调用渗入不相关的绘图操作。

\texttt{wipe()}函数替代了Processing的\texttt{background()}，使用了更简短、更主动的动词。这种命名是刻意的："wipe"（擦除）暗示一种物理动作（擦除表面），而非属性设置。

\subsection{无画布配置}

在Processing中，\texttt{setup()}的第一行通常是\texttt{size(w, h)}。在p5.js中，是\texttt{createCanvas(w, h)}。在AC中，没有等价的调用。运行时以设备原生分辨率提供画布，作品通过只读属性\texttt{screen.width}和\texttt{screen.height}接收。

这反映了移动优先的设计理念：在手机和平板上，"正确的"画布大小就是设备屏幕。允许程序员指定尺寸会在固定画布和可变视口之间产生不匹配。AC通过使画布默认自适应来消除这种不匹配。

% ============ 5. API SURFACE COMPARISON ============

\section{API接口比较}
\label{sec:comparison}

表~\ref{tab:lifecycle}比较了生命周期模型。表~\ref{tab:drawing}比较了绘图原语。

\begin{table}[h]
\small
\centering
\begin{tabularx}{\columnwidth}{lXXX}
\toprule
 & \textbf{Processing} & \textbf{p5.js} & \textbf{AC} \\
\midrule
初始化 & \texttt{setup()} & \texttt{setup()} & \texttt{boot(\$)} \\
渲染 & \texttt{draw()} & \texttt{draw()} & \texttt{paint(\$)} \\
逻辑 & （在draw中） & （在draw中） & \texttt{sim(\$)} \\
输入 & 回调 & 回调 & \texttt{act(\$)} \\
清理 & --- & \texttt{remove()} & \texttt{leave(\$)} \\
\midrule
渲染频率 & 60fps & 60fps & 显示Hz \\
逻辑频率 & 60fps & 60fps & 120fps固定 \\
作用域 & 全局 & 全局/实例 & 注入 \\
\bottomrule
\end{tabularx}
\caption{生命周期比较。\texttt{\$}表示解构式API注入。}
\label{tab:lifecycle}
\end{table}

\begin{table}[h]
\small
\centering
\begin{tabularx}{\columnwidth}{lXX}
\toprule
\textbf{Processing/p5} & \textbf{AC} & \textbf{说明} \\
\midrule
\texttt{background()} & \texttt{wipe()} & 动词，非属性 \\
\texttt{fill() + stroke()} & \texttt{ink()} & 统一、可链式 \\
\texttt{ellipse()} & \texttt{circle()} & 按形状命名 \\
\texttt{rect()} & \texttt{box()} & 更短名称 \\
\texttt{line()} & \texttt{line()} & 不变 \\
\texttt{point()} & \texttt{plot()} & 图形隐喻 \\
\texttt{text()} & \texttt{write()} & 主动动词 \\
--- & \texttt{flood()} & 区域填充 \\
--- & \texttt{paste()} & 缓冲区合成 \\
--- & \texttt{form()} & 3D渲染 \\
\bottomrule
\end{tabularx}
\caption{绘图原语比较。}
\label{tab:drawing}
\end{table}

\subsection{命名哲学}

Processing的API名称是描述性的名词和形容词：\texttt{background}、\texttt{ellipse}、\texttt{rect}、\texttt{stroke}、\texttt{fill}。它们描述正在设置或绘制\emph{什么}。AC的名称是主动动词：\texttt{wipe}、\texttt{ink}、\texttt{plot}、\texttt{write}、\texttt{flood}、\texttt{paste}。它们描述\emph{程序员正在做什么}——对表面执行的物理动作。

从描述性到祈使性命名的转变与乐器隐喻一致：乐器的界面通过动词来理解（按、吹、击、拨），而非名词。AC API读起来像一系列动作："擦除屏幕，上墨为粉色，画一个圆，写一些文字。"

\subsection{颜色模型}

Processing使用独立的\texttt{fill()}和\texttt{stroke()}调用，每个接受RGB、HSB或十六进制值。填充/描边的区分映射到矢量图形（SVG、PostScript），其中形状具有内部颜色和边界颜色。

AC使用单一的\texttt{ink()}调用为所有后续操作设置颜色。没有填充/描边区分：\texttt{circle()}是填充还是描边取决于其参数，而非环境上下文。\texttt{ink()}函数接受：

\begin{itemize}
  \item RGB值：\texttt{ink(255, 100, 50)}
  \item 命名颜色：\texttt{ink("pink")}、\texttt{ink("navy")}
  \item 灰度：\texttt{ink(128)}
  \item Alpha：\texttt{ink(255, 0, 0, 128)}
\end{itemize}

命名颜色对初学者来说是一个重要的可用性选择。在Processing需要\texttt{fill(255, 192, 203)}来获得粉色的地方，AC接受\texttt{ink("pink")}。命名颜色词汇表刻意保持小巧且富有唤起力，更接近颜料盒而非调色板。

% ============ 6. STATE MANAGEMENT ============

\section{状态管理}
\label{sec:state}

\subsection{Processing：全局变量}

\begin{lstlisting}[style=processingstyle]
int x = 100;
int y = 100;

void draw() {
  background(0);
  x += 1;
  ellipse(x, y, 20, 20);
}
\end{lstlisting}

Processing中的状态存在于全局变量中。这很简单，但有众所周知的问题：所有状态在任何地方都可变，没有封装，程序状态与API状态不可区分。

\subsection{AC：模块级闭包}

\begin{lstlisting}[style=acjsstyle]
let x = 100;
let y = 100;

function sim() {
  x += 1;
}

function paint({ wipe, ink, circle }) {
  wipe(0);
  ink(255).circle(x, y, 20);
}

export { sim, paint };
\end{lstlisting}

AC作品是ES模块。在模块作用域声明的变量对作品是私有的——对运行时、其他作品和全局作用域不可见。\texttt{export}语句只使生命周期函数可见。这不是文档强制执行的惯例；而是JavaScript模块系统的保证。

状态修改（\texttt{sim}）与状态渲染（\texttt{paint}）的分离鼓励了一种纪律，其中\texttt{paint}是模块状态的纯函数——它读取变量并绘制，但不修改它们。这不是强制的，但API结构使其成为自然模式。

% ============ 7. THE GENEALOGY ============

\section{谱系}
\label{sec:genealogy}

\subsection{Design By Numbers（1999）}

Maeda的Design By Numbers~\citep{maeda2001dbn}引入了核心洞察：一个用于视觉输出的编程环境，具有最简单的API。DBN没有动画循环——程序从上到下运行一次。原语集是最小的：\texttt{Paper}（背景）、\texttt{Pen}（线条粗细）、\texttt{Line}、\texttt{Set}（像素）和控制流。整个语言可以在一个下午学完。

AC继承了DBN的命名极简主义（简短的主动动词）及其确信API应该是可以穷尽学习的。

\subsection{Processing（2001）}

Processing~\citep{reas2003processing}添加了DBN所缺少的：动画（\texttt{draw()}循环）、交互（\texttt{mouseX/Y}）和更丰富的原语集。它还添加了DBN刻意避免的：状态（图形上下文）、标准库（数学、排版、图像加载）和编译步骤（Java）。Processing最伟大的创新不是任何单一函数，而是\texttt{setup()}/\texttt{draw()}这对组合：交互动画最简单的表达。

AC继承了生命周期函数模型，但将\texttt{draw()}拆分为三个关注点（\texttt{paint}、\texttt{sim}、\texttt{act}）。

\subsection{p5.js（2014）}

p5.js~\citep{mccarthy2015p5js}将Processing带入浏览器，使草图通过URL即时可分享。这是AC的前提条件：浏览器作为创意平台。p5.js还引入了实例模式，预示了AC的API注入模式。

AC继承了浏览器原生假设和可分享性原则（每个作品都是一个URL），但用完整运行时替代了页面中的库模型：作品不包含\texttt{<script>}标签；它们由平台加载。

\subsection{openFrameworks（2005）}

openFrameworks~\citep{openframeworks2005}使用了类似的生命周期（\texttt{setup()}、\texttt{update()}、\texttt{draw()}），但使用C++，将更新逻辑与渲染分开。AC的\texttt{sim()}/\texttt{paint()}分离更接近这个模型，而非Processing的统一\texttt{draw()}。

\subsection{乐器转向}

谱系可以总结为一系列分离：

\begin{enumerate}
  \item \textbf{DBN}：一次性执行（无动画循环）。
  \item \textbf{Processing}：初始化 + 渲染（\texttt{setup} + \texttt{draw}）。
  \item \textbf{openFrameworks}：初始化 + 更新 + 渲染。
  \item \textbf{AC}：初始化 + 输入 + 更新 + 渲染 + 清理（\texttt{boot} + \texttt{act} + \texttt{sim} + \texttt{paint} + \texttt{leave}）。
\end{enumerate}

每一步分离了先前混淆的关注点。AC的五函数模型是对交互程序生命周期最完整的分解，将其归入命名的、单一职责的入口点。结果类似于具身认知的感知-行动循环~\citep{varela1991embodied}：感知（\texttt{act}）、思考（\texttt{sim}）、行动（\texttt{paint}），以及出生（\texttt{boot}）和死亡（\texttt{leave}）两个额外的端点。

% ============ 8. WORKED EXAMPLE ============

\section{实例演示：弹球}
\label{sec:example}

为使比较具体化，我们在三个系统中实现一个弹球。

\subsection{Processing}

\begin{lstlisting}[style=processingstyle]
float x = 200, y = 200;
float vx = 3, vy = 2;

void setup() {
  size(400, 400);
}

void draw() {
  background(0);
  x += vx; y += vy;
  if (x < 0 || x > width) vx *= -1;
  if (y < 0 || y > height) vy *= -1;
  fill(255, 100, 180);
  noStroke();
  ellipse(x, y, 30, 30);
}
\end{lstlisting}

逻辑和渲染交织在\texttt{draw()}中。状态修改（\texttt{x += vx}）、边界检查和绘制共享一个函数。

\subsection{p5.js}

\begin{lstlisting}[style=p5style]
let x = 200, y = 200;
let vx = 3, vy = 2;

function setup() {
  createCanvas(400, 400);
}

function draw() {
  background(0);
  x += vx; y += vy;
  if (x < 0 || x > width) vx *= -1;
  if (y < 0 || y > height) vy *= -1;
  fill(255, 100, 180);
  noStroke();
  ellipse(x, y, 30, 30);
}
\end{lstlisting}

结构与Processing相同。唯一的变化是语法上的：\texttt{let}对比\texttt{float}，\texttt{createCanvas}对比\texttt{size}。

\subsection{AC}

\begin{lstlisting}[style=acjsstyle]
let x, y, vx = 3, vy = 2;

function boot({ screen }) {
  x = screen.width / 2;
  y = screen.height / 2;
}

function sim({ screen }) {
  x += vx; y += vy;
  if (x < 0 || x > screen.width) vx *= -1;
  if (y < 0 || y > screen.height) vy *= -1;
}

function paint({ wipe, ink, circle }) {
  wipe(0);
  ink(255, 100, 180).circle(x, y, 15);
}

export { boot, sim, paint };
\end{lstlisting}

关注点被分离：\texttt{boot}根据实际屏幕初始化状态，\texttt{sim}以120Hz更新物理（确保跨显示刷新率的一致行为），\texttt{paint}渲染而不修改状态。\texttt{ink().circle()}链式调用替代了三行代码（\texttt{fill}、\texttt{noStroke}、\texttt{ellipse}）。

% ============ 9. DESIGN RATIONALE ============

\section{设计理由}
\label{sec:rationale}

\subsection{为什么是五个函数？}

五函数模型不是出于软件工程的严谨性（作为教条的"关注点分离"），而是出于乐器隐喻。乐器有不同的阶段：你拿起它（boot），你倾听和回应（act），你思考接下来做什么（sim），你演奏（paint），最终你放下它（leave）。这些不是可以互换的时刻——你不会在演奏时调音，也不会在收起乐器时演奏。生命周期函数将这些区别形式化。

\subsection{为什么注入API？}

全局API对初学者方便，但创建了对平台的隐式依赖。如果不知道\texttt{ellipse()}是Processing函数而非用户函数，就无法理解Processing草图。AC的解构模式使依赖显式化：如果\texttt{circle}出现在函数体中，它就出现在函数签名中。这与ES模块导入背后相同的设计原则一致，应用在函数层面。

实际好处是热重载：因为作品不捕获全局引用，运行时可以在重载之间替换API对象而不会使作品的闭包失效。

\subsection{为什么是120Hz模拟？}

Processing的\texttt{draw()}以显示频率（通常60fps）运行。这意味着物理模拟、动画时序和输入处理都以与像素输出相同的频率运行。在144Hz显示器上，Processing草图运行速度快2.4倍；在30Hz显示器上，运行速度减半。

AC通过以固定120Hz运行\texttt{sim()}来解耦模拟与显示。这确保了确定性行为：弹球以相同的速度移动，无论显示刷新率如何。120Hz的频率被选为常见显示频率（60、120）的倍数，以及输入响应性的合理上限。

\subsection{为什么没有画布大小？}

要求\texttt{size()}或\texttt{createCanvas()}假设程序员知道（并关心）输出分辨率。在桌面上，这是合理的。在移动端——AC的主要目标——这是一个摩擦源：``正确的''大小取决于设备、方向和像素密度，所有这些程序员都不应该需要管理。AC将屏幕作为既定事实提供，就像一块已经绷好和涂底的画布。

% ============ 10. ADDITIONAL LIFECYCLE ============

\section{扩展生命周期}
\label{sec:extended}

除了核心五个函数外，AC作品还可以导出额外的生命周期钩子：

\begin{itemize}
  \item \texttt{beat(\$api)} --- 每个节拍器tick调用一次。BPM通过\texttt{sound.bpm}设置。在Processing中没有等价物；它反映了AC对音乐应用的取向。
  \item \texttt{preview(\$api)} --- 渲染作品的静态缩略图。用于作品列表和社交分享。
  \item \texttt{icon(\$api)} --- 渲染浏览器标签页的favicon。
  \item \texttt{meta()} --- 返回\texttt{\{title, desc\}}元数据用于Open Graph标签。
  \item \texttt{receive(event)} --- 处理来自其他作品或系统的消息，实现作品间通信。
\end{itemize}

\texttt{beat()}函数尤其值得注意。它将节奏时序从程序员必须实现的功能（通过Processing中的\texttt{millis()}）提升为第一等生命周期事件。这是API层面对音乐和节奏是核心用例而非附加功能的声明。

% ============ 11. CONCLUSION ============

\section{结论}
\label{sec:conclusion}

Processing的核心洞察——面向创意工作的编程环境应该提供生命周期，而非仅仅提供语言——存在于每个AC作品的核心。从\texttt{setup()}/\texttt{draw()}到\texttt{boot()}/\texttt{paint()}/\texttt{act()}/\texttt{sim()}/\texttt{leave()}的路径不是对Processing的背离，而是关于Processing模型在进一步扩展时能力边界的论证。

Processing证明了\texttt{setup()} + \texttt{draw()}足以使动画变得可及。AC作品API论证了相同的即时模式、生命周期驱动的方法——当分解为五个函数而非两个时——不仅可以支持速写，还可以支持\emph{练习}：反复回到同一作品，完善它，用它表演，随时间建立流畅度。

速写本隐喻（编写、运行、丢弃）和乐器隐喻（启动、演奏、练习、回归）并不对立。它们是同一连续体上的点。Processing展示了该连续体的存在。AC论证它比我们想象的延伸得更远——\texttt{setup()} + \texttt{draw()}是一个理念的开端，其完整形态需要\texttt{boot} + \texttt{act} + \texttt{sim} + \texttt{paint} + \texttt{leave}。

API即乐器。Processing是其核心。API如何分解编程行为决定了能创造什么。

\vspace{0.5em}
\noindent\textit{翻译自英文原版。原版请访问 \url{https://papers.aesthetic.computer}}

\bibliographystyle{plainnat}
\bibliography{references}

\end{document}
